\appendixitem{REPRESENTATIONS OF COMPACT GROUPS}

Let \(G\) be a compact topological group, and let \(\rho : G \to \mathbf{GL}(W)\) be a continuous representation of \(G\) on a finite dimensional complex vector space. We shall say that \(\rho\) is irreducible of real, complex or quaternionic type if this is the case for the \(\mathbf{C}^{(G)}\) -module \(W\) (relative to the algebra \(A = \mathbf{R}^{(G)}\) ). Let \(H\) be a separating positive hermitian form on \(W\) , invariant under \(G\) .

PROPOSITION 3. Assume that \(\rho\) is irreducible.

\begin{enumerate}
\def\labelenumi{\alph{enumi})}
\item
  The representation \(\rho\) is of real type if and only if there exists a non-zero symmetric bilinear form B on W, invariant under G. In this case the form B is separating; the set V of \(w \in W\) such that \(\mathrm{H}(w, x) = \mathrm{B}(w, x)\) for all \(x \in W\) is an \(\mathbf{R}\)-structure on W invariant under G.
\item
  The representation \(\rho\) is of complex type if and only if there exists no non-zero bilinear form on W invariant under G.
\item
  The representation \(\rho\) is of quaternionic type if and only if there exists a non-zero alternating bilinear form on W, invariant under G; such a form is necessarily separating.
\end{enumerate}

For \(\theta \in \mathrm{Hom}_{\mathbf{C}^{(\mathrm{G})}}(\mathrm{W},\overline{\mathrm{W}})\) and \(x,y\in \mathrm{W}\) , put \(\mathrm{B}_{\theta}(x,y) = \mathrm{H}(\theta x,y)\) . Then \(\mathrm{B}_{\theta}\) is a bilinear form on \(\mathrm{W}\) , invariant under \(\mathrm{G}\) , and separating if \(\theta\) is non-zero. Denote by \(\mathcal{B}(\mathrm{W})^{\mathrm{G}}\) the space of bilinear forms on \(\mathrm{W}\) invariant under \(\mathrm{G}\) ; the map \(\theta \mapsto \mathrm{B}_{\theta}\) from \(\mathrm{Hom}_{\mathbf{C}^{(\mathrm{G})}}(\mathrm{W},\overline{\mathrm{W}})\) to \(\mathcal{B}(\mathrm{W})^{\mathrm{G}}\) is an isomorphism of \(\mathbf{C}\) -vector spaces. This implies, in particular, assertion \(b\) ).

Let \(\theta\) be a \(\mathbf{C}^{(\mathrm{G})}\) -isomorphism from W to \(\overline{\mathrm{W}}\) such that \(\theta \circ \theta = \alpha_{\mathrm{W}}\) , with \(\alpha \in \{-1, +1\}\) (Prop. 2); since \(\mathcal{B}(\mathrm{W})^{\mathrm{G}}\) is of dimension 1, there exists \(\varepsilon \in \mathbf{C}\) such that

\[
\mathrm{B} _ {\theta} (y, x) = \varepsilon \mathrm{B} _ {\theta} (x, y) \quad \text { for   all } x, y \text { in } \mathrm{W}.
\]

Iterating, we obtain \(\mathrm{B}_{\theta}(y,x)=\varepsilon\mathrm{B}_{\theta}(x,y)=\varepsilon^{2}\mathrm{B}_{\theta}(y,x)\) , so \(\varepsilon^{2}=1\) and \(\varepsilon\in\{-1,+1\}\) . Moreover, for x in W,

\[
\mathrm{H} (\theta x, \theta x) = \mathrm{B} _ {\theta} (x, \theta x) = \varepsilon \mathrm{B} _ {\theta} (\theta x, x) = \varepsilon \mathrm{H} (\theta \circ \theta (x), x) = \varepsilon \alpha \mathrm{H} (x, x)
\]

so \(\varepsilon\alpha > 0\) since H is positive, that is, \(\varepsilon = \alpha\) . Assertions a) and c) now follow from Prop. 2.

Denote by dg the Haar measure of total mass 1 on G.

Lemma 1. Let \(W^{G}\) be the subspace of W consisting of the elements invariant under G. The endomorphism \(\int_{G}\rho(g)dg\) of W is a projection with image \(W^{G}\) , compatible with the operations of G. In particular,

\[
\dim \mathrm{W} ^ {\mathrm{G}} = \int_ {\mathrm{G}} \operatorname{Tr} \rho (g) d g.
\]

Put \(p = \int_{\mathrm{G}} \rho(g) dg\) ; for \(h \in \mathrm{G}\) ,

\[
\rho (h) \circ p = \int_ {\mathrm{G}} \rho (h g) d g = \int_ {\mathrm{G}} \rho (g) d g = p
\]

and similarly \(p \circ \rho(h) = p\) . Thus, p is compatible with the operations of G, and its image is contained in \(W^{G}\) . If \(w \in W^{G}\) , we have \(p(w) = \int_{\mathbf{G}} \rho(g) w \, dg = w\) , hence the lemma.

Lemma 2. Let \(u\) be an endomorphism of a finite dimensional vector space \(E\) over a field \(K\) . Then

\[
\operatorname{Tr} u ^ {2} = \operatorname{Tr} \mathbf {S} ^ {2} (u) - \operatorname{Tr} \Lambda^ {2} (u).
\]

Let \(\chi_u(\mathrm{X}) = \prod_{i=1}^{n} (\mathrm{X} - \alpha_i)\) be a decomposition of the characteristic polynomial of \(u\) into linear factors in a suitable extension of K. We have \(\operatorname{Tr} u^2 = \sum_i \alpha_i^2\) , \(\operatorname{Tr} \Lambda^2(u) = \sum_{i < j} \alpha_i \alpha_j\) , \(\operatorname{Tr} \mathbf{S}^2(u) = \sum_{i \leq j} \alpha_i \alpha_j\) (cf.~Algebra, Chap. VII, §5, no. 5, Cor. 3), hence the result.

PROPOSITION 4. Assume that \(\rho\) is irreducible. Then, \(\rho\) is of real (resp. complex, resp. quaternionic) type if and only if the integral \(\int_{\mathrm{G}} \mathrm{Tr} \rho(g^2) dg\) is equal to 1 (resp. 0, resp. -1).

Denote by \(\check{\rho}\) the contragredient representation of \(\rho\) on \(W^{*}\) (defined by \(\check{\rho}(g)={}^{t}\rho(g^{-1})\) ). Applying Lemma 2 to \(\check{\rho}(g)\) and integrating over G gives

\[
\int_ {\mathrm{G}} \operatorname{Tr} \rho (g ^ {2}) d g = \int_ {\mathrm{G}} \operatorname{Tr} ^ {t} \rho (g ^ {- 2}) d g = \int_ {\mathrm{G}} \operatorname{Tr} \mathbf {S} ^ {2} (\check {\rho} (g)) d g - \int_ {\mathrm{G}} \operatorname{Tr} \Lambda^ {2} (\check {\rho} (g)) d g
\]

hence, by Lemma 1,

\[
\int_ {\mathrm{G}} \mathrm{Tr} \rho (g ^ {2}) d g = \dim (\mathbf {S} ^ {2} \mathrm{W} ^ {*}) ^ {\mathrm{G}} - \dim (\Lambda^ {2} \mathrm{W} ^ {*}) ^ {\mathrm{G}}.
\]

But \(S^{2}W^{*}\) (resp. \(\Lambda^{2}W^{*}\) ) can be identified with the space of symmetric (resp. alternating) bilinear forms on W. Thus, the proposition follows immediately from Prop. 3.

